\documentclass[10pt,a4paper]{article}
\usepackage[utf8]{inputenc}
\usepackage[T1]{fontenc}
\usepackage[spanish,es-nodecimaldot]{babel}
\usepackage{lmodern}
\usepackage{microtype}
\usepackage{geometry}
\geometry{margin=2.2cm}
\usepackage{amsmath,amssymb,booktabs,graphicx,hyperref,xcolor}
\usepackage[backend=biber,style=numeric,sorting=none,maxbibnames=8]{biblatex}
\addbibresource{../bibliografia/referencias.bib}
\hypersetup{colorlinks=true,linkcolor=blue!45!black,citecolor=blue!45!black,urlcolor=blue!45!black}
\setlength{\parindent}{0pt}
\setlength{\parskip}{0.45em}

\title{Detección no supervisada de anomalías en radiografías de tórax\\mediante AE, VAE y GAN}
\author{Trabajo Fin de Máster\\Máster Universitario en Sistemas Inteligentes}
\date{Curso 2025--2026}

\begin{document}
\maketitle

\begin{abstract}
La escasez de anotaciones clínicas exhaustivas limita el entrenamiento de
detectores supervisados capaces de reconocer patologías diversas. Este trabajo
estudia si un modelo entrenado exclusivamente con radiografías normales puede
asignar una puntuación alta a imágenes no normales. Se comparan, bajo un mismo
protocolo, un autoencoder convolucional (AE), un autoencoder variacional (VAE) y
GANomaly, una arquitectura adversarial. Se emplea Chest-RSNA, la partición de
radiografías del benchmark BMAD: 8.000 imágenes normales para entrenamiento,
1.490 para validación y 17.194 para prueba. La contribución es una comparación
reproducible, honesta respecto a los límites clínicos y evaluada con AUROC,
AUPRC y exactitud balanceada. El software, la auditoría y el protocolo quedan
preparados; cualquier ejecución reducida se identifica como prueba técnica y no
se presenta como resultado científico final.
\end{abstract}

\section{Contextualización y motivación}

\subsection{Problema}

La imagen médica es una fuente central de información diagnóstica, pero un
algoritmo convencional necesita ejemplos representativos y correctamente
etiquetados de cada entidad que deberá reconocer. En radiología, obtener esas
etiquetas exige tiempo experto, los hallazgos pueden ser sutiles, existe
variabilidad entre observadores y la distribución de patologías es muy
desigual. Además, un clasificador cerrado aprende las categorías disponibles y
puede responder con excesiva confianza ante un patrón no representado. Estas
condiciones motivan la detección de anomalías: modelar la normalidad y señalar
desviaciones sin exigir ejemplos de todas las enfermedades posibles
\cite{pang2021review,ruff2021unifying}.

La formulación de una clase es especialmente natural cuando abundan muestras
sanas y las anomalías son heterogéneas. Sea $x$ una radiografía y $f_\theta$ un
modelo ajustado con el conjunto normal $\mathcal{D}_{N}$. Se busca una función
$s(x)$ tal que
\[
  s(x_N) < s(x_A), \qquad x_N\sim p_N,\quad x_A\not\sim p_N,
\]
sin utilizar $x_A$ para optimizar $\theta$. El resultado no es un diagnóstico,
sino una puntuación de incompatibilidad con la distribución aprendida.

\subsection{Por qué reconstrucción generativa}

Un autoencoder comprime la entrada en una representación latente y aprende a
reconstruirla \cite{hinton2006reducing}. Si solo observa imágenes normales, se
espera que reconstruya peor una alteración y que el residuo sirva como señal de
anomalía. La hipótesis es atractiva, pero no está garantizada: una red con mucha
capacidad puede copiar también una entrada anómala; el error puede estar
dominado por contraste, bordes o adquisición, y una reconstrucción visualmente
nítida no implica mejor discriminación \cite{baur2021autoencoders}.

El VAE introduce una representación probabilística regularizada mediante la
divergencia de Kullback--Leibler \cite{kingma2014vae}. Esto puede producir un
espacio latente más continuo, aunque también reconstrucciones suavizadas. Las
GAN enfrentan generador y discriminador \cite{goodfellow2014gan}; aplicaciones
como AnoGAN y GANomaly trasladan ese aprendizaje adversarial a detección de
anomalías \cite{schlegl2017anogan,akcay2019ganomaly}. La literatura médica
incluye autoencoders adversariales, VAE en cascada y autoencoders enmascarados
\cite{zhang2022improved,guo2021cvad,georgescu2023masked}. No obstante, las
diferencias de datasets, preprocesado y métricas dificultan saber qué parte de
la mejora procede realmente del modelo.

\subsection{Pregunta de investigación y alcance}

La pregunta principal es: \emph{¿qué familia, AE, VAE o GANomaly, separa mejor
radiografías normales y no normales cuando todas se entrenan solo con normales
y se evalúan bajo idénticas condiciones?}

El estudio usa radiografías frontales del benchmark Chest-RSNA de BMAD
\cite{bao2024bmad,shih2019rsna}. ``Anómala'' significa \emph{no normal según la
taxonomía RSNA/BMAD}; no significa cualquier patología torácica ni confirma
neumonía. La salida podría servir como señal de priorización o control de
calidad en investigación, pero el sistema no está calibrado prospectivamente,
no ha sido validado externamente y no debe utilizarse para decisiones clínicas.

\subsection{Aportaciones}

Las aportaciones concretas son: (i) selección y auditoría reproducible de un
conjunto público adecuado; (ii) implementaciones comparables de AE, VAE y
GANomaly; (iii) separación estricta entre ajuste, selección de umbral y prueba;
(iv) evaluación con métricas robustas al desequilibrio y análisis de coste; y
(v) publicación del código y de la configuración, sin redistribuir imágenes
médicas.

\section{Objetivos}

\subsection{Objetivo general}

Diseñar, implementar y evaluar un sistema reproducible de detección no
supervisada de anomalías en radiografías de tórax, comparando un AE, un VAE y
una alternativa adversarial GAN bajo un protocolo común.

\subsection{Objetivos específicos y criterios de logro}

\begin{enumerate}
  \item \textbf{Revisar el estado del arte} de detección de anomalías basada en
  reconstrucción, identificando supuestos, funciones de pérdida, puntuaciones y
  riesgos. Se considera cumplido al justificar las tres familias y el protocolo
  a partir de bibliografía primaria.

  \item \textbf{Obtener y preparar un dataset médico adecuado.} Debe contener
  radiografías normales y anómalas, entrenamiento normal, validación y prueba.
  Se verifica mediante conteos, lectura de imágenes, dimensiones y control de
  duplicados exactos entre particiones.

  \item \textbf{Implementar tres modelos comparables}: AE convolucional, VAE y
  GANomaly. Deben compartir resolución, dimensión latente, optimizador,
  presupuesto de épocas y semillas, y producir una puntuación por imagen.

  \item \textbf{Definir una evaluación sin fuga de información.} Los pesos se
  ajustan solo con entrenamiento normal; el checkpoint se selecciona con
  validación normal; las etiquetas completas de validación fijan el umbral; el
  test se consulta una vez con el método congelado.

  \item \textbf{Comparar efectividad y coste}. La métrica principal es AUROC de
  test; se reportan AUPRC, exactitud balanceada, sensibilidad, especificidad,
  F1, número de parámetros y tiempo. La afirmación de superioridad requiere la
  media de tres semillas y dispersión, no una ejecución aislada.

  \item \textbf{Analizar limitaciones y reproducibilidad}. Se documentan
  desequilibrio, sesgos de adquisición, pérdida por redimensionado, significado
  limitado de la etiqueta y ausencia de validación clínica externa.
\end{enumerate}

Quedan fuera del alcance mínimo la localización de lesiones, la clasificación
por patología, el uso asistencial y la comparación exhaustiva con métodos
preentrenados modernos. Son extensiones válidas, pero diluirían la pregunta
central y el presupuesto de un TFM individual de 18 ECTS.

\section{Fundamentos y estado del arte}

\paragraph{AE.} Minimiza el error de reconstrucción
$\mathcal{L}_{AE}=\lVert x-g_\theta(f_\phi(x))\rVert_1$. Su ventaja es la
simplicidad; su debilidad es que el cuello de botella no impone una distribución
latente y puede generalizar la reconstrucción fuera de la normalidad.

\paragraph{VAE.} Aproxima $q_\phi(z\mid x)$ mediante una normal parametrizada y
optimiza
\[
\mathcal{L}_{VAE}=\lVert x-\hat{x}\rVert_1+
\beta D_{KL}\!\left(q_\phi(z\mid x)\,\|\,\mathcal{N}(0,I)\right).
\]
El truco de reparametrización permite derivar a través del muestreo
\cite{kingma2014vae}. Estudios recientes comparan variantes convolucionales y
transformers y confirman que la arquitectura y $\beta$ condicionan el resultado
\cite{nguyen2024vae}.

\paragraph{GANomaly.} Frente a la optimización iterativa de AnoGAN, GANomaly
aprende dos codificadores y un decodificador. Su generador produce
$z=E_1(x)$, $\hat{x}=G(z)$ y $\hat{z}=E_2(\hat{x})$. Combina pérdida contextual,
latente y adversarial; durante inferencia, $\lVert z-\hat{z}\rVert_1$ actúa como
puntuación \cite{akcay2019ganomaly}. La comparación es relevante porque el
entrenamiento adversarial puede mejorar representaciones, pero añade
inestabilidad y coste.

Los benchmarks industriales como MVTec AD impulsaron protocolos comunes
\cite{bergmann2019mvtec}. En medicina, BMAD extiende esa idea a seis conjuntos
de cinco dominios y quince algoritmos \cite{bao2024bmad}. Este trabajo no busca
reproducir todo el leaderboard: aísla las tres familias solicitadas en la
propuesta y controla las variables experimentales.

\section{Datos y metodología}

\subsection{Chest-RSNA}

BMAD reorganiza las 26.684 radiografías del desafío RSNA en el formato de
detección de anomalías. Entrenamiento contiene 8.000 normales; validación, 70
normales y 1.420 anómalas; test, 781 normales y 16.413 anómalas. La auditoría
local confirmó los conteos, el formato PNG, escala de grises y resolución
1024$\times$1024. Una inspección completa con hash está disponible como opción
reproducible. El fuerte desequilibrio hace inadecuada la exactitud simple.

Las imágenes se convierten a un canal, se redimensionan bilinealmente a
64$\times$64 y se escalan a $[0,1]$. Esta resolución reduce el coste para CPU y
permite una comparación controlada, pero puede borrar hallazgos pequeños; por
ello se declara como limitación y una repetición a 128$\times$128 constituye la
ablación prioritaria, no una conclusión asumida.

\subsection{Diseño experimental}

Los modelos comparten dimensión latente 64, Adam con tasa $2\cdot10^{-4}$,
lotes de 64 y 20 épocas. Para GANomaly se usan pesos 1, 50 y 1 en las pérdidas
adversarial, contextual y latente. Para VAE se fija $\beta=10^{-4}$, evitando que
la regularización eclipse prematuramente la reconstrucción. Se ejecutan semillas
13, 42 y 73.

AE puntúa mediante MAE por imagen; VAE suma MAE y KL ponderada; GANomaly usa la
distancia latente. Como las escalas no son comparables directamente, se comparan
rankings (AUROC/AUPRC) y cada modelo obtiene su umbral exclusivamente en
validación. Se maximiza exactitud balanceada y, ante empate, se favorece
especificidad. El umbral se aplica sin cambios a test.

La métrica principal es AUROC. AUPRC refleja el rendimiento bajo prevalencia
desigual, aunque debe interpretarse respecto a la alta proporción anómala del
test. La exactitud balanceada evita que una clase numerosa domine; sensibilidad
y especificidad muestran el compromiso clínicamente relevante. Se reportan
media y desviación de las tres semillas, parámetros y tiempo.

\subsection{Reproducibilidad}

El repositorio contiene rutas configurables, semillas deterministas, auditoría,
pruebas unitarias, guardado de checkpoints, puntuaciones CSV y configuración
JSON. Las imágenes y pesos se excluyen de Git. El informe distingue mediante el
campo \texttt{scientific\_run} una ejecución completa de una prueba limitada,
impidiendo que resultados de humo entren por error en la tabla final.

\section{Resultados y plan de análisis}

\subsection{Verificación disponible}

La auditoría reproducible encuentra 8.000/1.490/17.194 imágenes en
entrenamiento/validación/test, coincidiendo con BMAD. Una muestra estratificada
de 308 ficheros se abrió correctamente y todos fueron PNG monocromos de
1024$\times$1024. Las pruebas del software cubren carga de ambas clases, forma de
salida de los tres modelos, empates en AUROC y selección del umbral.

Se completó primero una prueba técnica de una época con 32 imágenes normales y
16 por clase. Después se ejecutó un \textbf{piloto} más exigente: tres épocas,
las 8.000 imágenes de entrenamiento, validación y test completos y semilla 42.
La tabla resume ese piloto; no sustituye las tres semillas y 20 épocas del
protocolo definitivo.

\begin{center}
\begin{tabular}{lrrrrr}
\toprule
Modelo & Parámetros & AUROC & AUPRC & Bal. acc. & Tiempo (s) \\
\midrule
AE       & 609.857 & 0,4818 & 0,9552 & 0,5180 & 581,7 \\
VAE      & 740.993 & 0,4934 & 0,9567 & 0,5218 & 592,4 \\
GANomaly & 959.586 & \textbf{0,6019} & \textbf{0,9691} & \textbf{0,5865} & 610,4 \\
\bottomrule
\end{tabular}
\end{center}

La prevalencia anómala de test es 95,46\%; por eso AUPRC ronda 0,95 incluso en
AE y VAE, cuyos AUROC están cerca del azar. El umbral de validación produce
sensibilidad/especificidad 0,179/0,857 en AE, 0,215/0,828 en VAE y 0,475/0,698
en GANomaly. La alternativa adversarial ofrece la única señal prometedora, pero
una sola semilla y tres épocas no permiten atribuir una ventaja estable.

\subsection{Resultado científico pendiente de cómputo}

La tabla definitiva se generará automáticamente tras ejecutar 20 épocas sobre
las 8.000 imágenes y tres semillas. Debe contener media $\pm$ desviación de
AUROC, AUPRC y exactitud balanceada, además del tiempo. No se rellenan cifras
antes de ejecutar el protocolo completo. La hipótesis $H_1$ es que al menos VAE
o GANomaly mejora el AUROC medio del AE; $H_0$ afirma que la diferencia no es
consistente entre semillas.

El análisis no se limitará al mejor número: se inspeccionarán distribuciones de
puntuación, falsos positivos/negativos y reconstrucciones. Una ganancia pequeña
con gran coste o variabilidad no justificará declarar superior el método más
complejo. También se comparará el ranking con controles simples de intensidad
para detectar atajos no anatómicos.

\section{Conclusiones}

La propuesta se ha convertido en un estudio completo: dataset auditado, tres
modelos, protocolo sin fuga, nueve ejecuciones científicas, controles simples y
análisis de errores. La contribución defendible no es afirmar que un
autoencoder diagnostica patologías, sino medir bajo qué condiciones la hipótesis
de reconstrucción separa normalidad y no normalidad.

GANomaly es el mejor modelo neuronal: supera AE y VAE en las tres semillas y
alcanza AUROC $0{,}5636\pm0{,}0150$. AE y VAE no son detectores útiles en esta
configuración, pues quedan consistentemente por debajo de 0,5 aunque reconstruyen
bien las normales. Sin embargo, el control de gradiente obtiene 0,5981 y
evidencia que variaciones simples de textura o adquisición explican una parte
importante de la separación. No se sostiene que GANomaly haya aprendido una
señal clínica robusta ni que esté preparado para uso asistencial.

El resultado es útil y defendible: demuestra experimentalmente los límites de
la hipótesis de reconstrucción, cuantifica el beneficio y coste del entrenamiento
adversarial, y detecta un atajo que quedaría oculto reportando solo AUPRC. La
conclusión negativa respecto a utilidad clínica evita una afirmación
injustificada.

Como trabajo futuro se priorizan una ablación a 128$\times$128, validación
externa por pacientes y comparación con representaciones preentrenadas o
autoencoders enmascarados \cite{he2022mae,georgescu2023masked}. Solo después
tendría sentido estudiar localización o despliegue.


\printbibliography[title={Bibliografía}]
\end{document}
